\section{Limitações dos Testes}
\label{sec:limitacoes-testes}

Apesar de o benchmark permitir uma comparação objetiva entre Morpheus, Sunny e Rocket, algumas limitações devem ser consideradas. A primeira delas está relacionada ao ambiente de execução. Os relatórios não detalham completamente as configurações de hardware e software utilizadas em cada cenário. Informações como processador exato, quantidade de núcleos, memória disponível, sistema operacional, versões de runtime, uso de container e limites de CPU ou RAM não estão completamente especificadas.

Essa ausência limita a reprodutibilidade exata dos testes. Em benchmarks computacionais, pequenas diferenças no ambiente podem afetar métricas como tempo de execução, consumo de memória e uso de CPU. Por isso, os resultados devem ser entendidos como válidos para o protocolo analisado, e não como medidas absolutas independentes do ambiente.

Outra limitação está relacionada ao controle da aleatoriedade. Algoritmos genéticos dependem de processos estocásticos, como seleção, \textit{crossover}, mutação e geração da população inicial. Os relatórios não informam as sementes aleatórias utilizadas nas execuções. Dessa forma, embora as repetições reduzam parte da variabilidade, a reprodução exata dos resultados fica limitada.

Também não há detalhamento completo e homogêneo sobre a composição do erro médio. Como essa métrica não é decomposta por tipo de violação de forma suficiente para sustentar comparação qualitativa aprofundada, não é possível identificar, a partir da síntese apresentada, se os erros restantes correspondem a choques de professor, choques de turma, horários inválidos, demandas não alocadas ou outro tipo de inconsistência. Essa limitação impede uma avaliação mais profunda da qualidade acadêmica das grades geradas.

Além disso, as métricas de CPU devem ser interpretadas com cautela. A coleta foi baseada na telemetria do Docker para os containers monitorados, em que 100\% de CPU representa aproximadamente o uso integral de um núcleo lógico, e não necessariamente a ocupação de todo o host físico. Ainda assim, os relatórios não detalham completamente possíveis interferências externas, concorrência, aquecimento, janela exata de amostragem ou descarte de outliers. Essa cautela é especialmente relevante no caso do Rocket, pois a média de CPU registrada foi baixa, enquanto os picos foram elevados, o que pode refletir características da coleta por amostragem, execução concorrente concentrada em intervalos curtos ou particularidades do ambiente de medição.

Outra limitação decorre do fato de os relatórios do Rocket terem sido gerados separadamente dos relatórios de Morpheus e Sunny, embora com o mesmo perfil de \textit{payload} e os mesmos cenários declarados. Essa separação não invalida a comparação descritiva, mas recomenda cautela adicional, pois diferenças de momento de execução, estado do host, versão de container ou condições de carga podem afetar métricas operacionais.

Por fim, a análise apresentada é predominantemente descritiva. Não foram aplicados testes estatísticos inferenciais, intervalos de confiança ou análise formal de significância entre os três algoritmos. Assim, os resultados devem ser interpretados como evidências experimentais do benchmark executado, e não como comprovação estatística universal de superioridade de um algoritmo sobre o outro.
